% ---- pg-472 （书页 460） ----

\noindent\textbf{例 11.4}\quad (血友病数据的费希尔线性判别)

例 11.3 中介绍了关于检查 A 型血友病基因携带者的研究. 我们记得, 在同等代价和同等先验概率下, 线性判别函数为
\[
\hat{y}=\hat{\mathbf{a}}'\mathbf{x}=[\bar{\mathbf{x}}_1-\bar{\mathbf{x}}_2]'\mathbf{S}_{\mathrm{p}}^{-1}\mathbf{x}=37.61x_1-28.92x_2
\]
该线性判别函数就是费希尔线性函数, 它最大限度地分离两个总体, 因而样本中的最大分离度为
\[
\begin{aligned}
D^2&=(\bar{\mathbf{x}}_1-\bar{\mathbf{x}}_2)'\mathbf{S}_{\mathrm{p}}^{-1}(\bar{\mathbf{x}}_1-\bar{\mathbf{x}}_2)\\
&=[0.241\,8,\ -0.065\,2]\begin{bmatrix}131.158&-90.423\\-90.423&108.147\end{bmatrix}\begin{bmatrix}0.241\,8\\-0.065\,2\end{bmatrix}\\
&=10.98
\end{aligned}
\]
费希尔对分离问题的解法也可以用于对新的观测值作分类.

\begin{mvframe}[基于费希尔判别函数的分类法则\textsuperscript{\textcircled{1}}]
若
\phantomsection\label{eq:11-25}
\begin{equation}
\begin{aligned}
\hat{y}_0&=(\bar{\mathbf{x}}_1-\bar{\mathbf{x}}_2)'\mathbf{S}_{\mathrm{p}}^{-1}\mathbf{x}_0\\
&\geqslant\hat{m}=\frac{1}{2}(\bar{\mathbf{x}}_1-\bar{\mathbf{x}}_2)'\mathbf{S}_{\mathrm{p}}^{-1}(\bar{\mathbf{x}}_1+\bar{\mathbf{x}}_2)
\end{aligned}
\tag{11-25}
\end{equation}
或
\[
\hat{y}_0-\hat{m}\geqslant0
\]
将 $\mathbf{x}_0$ 分到 $\pi_1$. 若
\[
\hat{y}_0<\hat{m}
\]
或
\[
\hat{y}_0-\hat{m}<0
\]
将 $\mathbf{x}_0$ 分到 $\pi_2$.
\end{mvframe}

\noindent{\footnotesize \textcircled{1}\quad 我们必须有 $(n_1+n_2-2)\geqslant p$, 否则 $\mathbf{S}_{\mathrm{p}}$ 为奇异阵, 因而通常的逆矩阵 $\mathbf{S}_{\mathrm{p}}^{-1}$ 不存在.}

\hyperref[fig:11.5]{图 11.5} 是 $p=2$ 时对方法(11-23)的示意性说明. 图中所有的点均投影到沿 $\hat{\mathbf{a}}$ 方向的一条直线上, 而这个方向在不断变化, 直到两个样本被最大限度地分开为止.

\begin{mvfigure}
\mvinc{pg-472}{0.251,0.643,0.736,0.846}
\phantomsection\label{fig:11.5}\mvcap{图 11.5\quad 费希尔方法对两个总体的图形表示($p=2$)}
\end{mvfigure}
